
\clearpage
\section{Aklin and Mildenberger --- \textit{Prisoners of the Wrong Dilemma: Why Distributive Conflict, Not Collective Action, Characterizes the Politics of Climate Change}}

\subsection{Core argument}

Aklin and Mildenberger challenge the dominant interpretation of climate change politics as a global collective-action problem driven primarily by free-riding. The conventional account treats a stable climate as a global public good: every country benefits from mitigation, but each individual government has an incentive to let others bear the costs of reducing emissions. From this perspective, the central political problem is to create international institutions that monitor behavior, discourage defection, and make cooperation reciprocal. The Kyoto Protocol, and to a lesser extent the Paris Agreement, reflect this logic by emphasizing international commitments, transparency, monitoring, and compliance.

The authors argue that this framework has had enormous theoretical influence despite surprisingly weak empirical support. They distinguish between a \textit{strong collective-action} account, according to which free-riding is the dominant constraint on climate policy, and a \textit{weak collective-action} account, according to which free-riding is merely one among several relevant considerations. Their main claim is that the strong version is contradicted by important patterns of climate policy, while the weaker version often becomes observationally equivalent to alternative explanations. In particular, many of the same observations can be explained more parsimoniously through domestic distributive conflict.

The key distinction is between climate \textit{outcomes} and climate \textit{policy}. Climate outcomes do involve global externalities: emissions generated in one country contribute to a global problem. Climate policy, however, is made through political institutions and creates domestic winners and losers. Decarbonization changes the distribution of income, employment, investment, and political influence across sectors. Consequently, governments face conflict between actors that expect to benefit from climate reform and actors that expect to lose. The politics of climate change may therefore be better understood as a struggle over distribution within countries rather than primarily as a Prisoner's Dilemma among countries.

\subsection{Collective-action predictions and the distributive alternative}

The strong collective-action framework generates three central predictions. First, effective climate cooperation requires international institutions capable of monitoring behavior and punishing noncompliance. Second, countries should adopt costly climate policies reciprocally: unilateral action should be irrational when others are able to free-ride. Third, defection by pivotal states should induce defection or at least substantial retrenchment by others, because the absence of a major emitter sharply reduces the benefits of continued cooperation.

The weaker version relaxes these expectations. International negotiations should at least feature concerns about free-riding; other states' actions should increase the probability that a country adopts costly climate policy; and defection by major players should increase the probability that others reduce their own efforts. This weaker formulation is more difficult to falsify, but it also becomes less distinctive because many forms of interdependence can generate similar behavior without free-riding being the underlying mechanism.

The authors contrast these expectations with a \textbf{distributive-conflict framework}. Climate policy reorganizes economic activity and therefore divides political actors according to expected gains and losses. One family of distributive theories emphasizes the influence of organized special interests, particularly carbon-intensive sectors with close ties to policy makers. Another emphasizes the balance of power between ``green'' and ``brown'' coalitions, including firms, unions, voters, environmental organizations, parties, and other actors. Domestic institutions determine which interests receive greater political weight.

This framework yields a different set of predictions. Countries should act when pro-climate forces are strong enough to overcome or neutralize opponents, regardless of what foreign governments do. Foreign defection need not be reciprocated because policy depends primarily on domestic political configurations. Finally, meaningful climate policy can emerge even in the absence of international institutions designed to manage free-riding.

\subsection{Evidence from national climate policy}

The first major empirical challenge to the collective-action account is the existence of extensive unilateral climate policy. Countries have repeatedly adopted carbon pricing, renewable-energy support, and other mitigation measures without waiting for reciprocal action by other states.

The United States rejected the Kyoto Protocol under the George W. Bush administration, yet the EU did not respond by abandoning climate action. Instead, the EU developed the Emissions Trading Scheme.

\textbf{These developments are difficult to reconcile with the strongest collective-action predictions. If climate action were fundamentally reciprocal, the withdrawal of a pivotal emitter such as the United States should have reduced incentives} for others to proceed. Yet unilateral policies expanded.

\textbf{This pattern is directly compatible with the distributive-conflict view}: governments act when domestic coalitions favoring climate policy are politically strong enough, even if other countries do not act.

\subsection{Public opinion and conditional cooperation}

The authors next review research on mass attitudes. \textbf{A collective-action account implies that citizens should behave as conditional cooperators: support for costly domestic climate policy should depend heavily on whether other countries also contribute}. 

This is not true, emprirically it is true that large majorities in several countries support climate action even when other countries do not act. Some respondents even behave as ``counterbalancers'', becoming more supportive of domestic action when others fail to cooperate.

Some studies find that broader international participation increases support for climate agreements. However, Aklin and Mildenberger emphasize that this does not establish that free-riding is the causal mechanism. Conditionality can emerge from fairness, reciprocity, social norms, perceptions of legitimacy, or broader political cues. Evidence that people respond to what others do is therefore not the same as evidence that they are motivated by fears of being exploited as free-riders.

The distinction between \textit{conditional action} and \textit{conditional cooperation} is central.

\subsection{Political elites and the cases of Kyoto}

The authors revisit political episodes often cited as evidence of conditional cooperation. The 1997 Byrd--Hagel resolution in the US Senate rejected participation in a climate agreement that imposed obligations on industrialized countries while exempting major developing countries. At first glance, this appears to be a textbook reaction to free-riding. A closer examination, however, reveals that the coalition supporting the resolution was largely built around domestic distributive interests.

The same logic helps reinterpret George W. Bush's withdrawal from Kyoto. Although the administration publicly referred to the absence of binding commitments for China and India. \textbf{The actors pushing most strongly against climate action were not conditional cooperators who would have accepted regulation under a more reciprocal agreement}. They were largely unconditional opponents of costly domestic mitigation. Free-riding rhetoric helped legitimate their bargaining position, but it did not necessarily reveal the underlying cause of their preferences.

Canadian governments sometimes justified inaction by pointing to US policy, but political leaders who used this argument often opposed climate regulation even when US policy changed.

In Australia, ratification of Kyoto followed a change in government rather than a change in the international structure of incentives. What changed was the domestic balance of political power.

\subsection{Inference and conclusion}

These cases lead to two important methodological lessons.

The first is \textit{equifinality}: an outcome that is consistent with collective-action theory may also be consistent with distributive politics. Political inaction after another country defects does not prove that the second country was motivated by free-riding. Domestic opponents may have preferred inaction regardless and simply used foreign behavior as political justification.

The second problem is \textit{epiphenomenality}. Political rhetoric about reciprocity may be a strategic instrument rather than evidence of actors' true motivations. Carbon-intensive interests have incentives to frame their opposition in internationally legitimate terms, and free-riding arguments can be useful for this purpose.

The authors therefore conclude that the strongest collective-action model lacks empirical support, while weaker versions are difficult to distinguish from competing explanations. Their alternative is not to claim that international politics or free-riding never matter. Rather, the empirical burden should be reversed: researchers should demonstrate when and how free-riding is actually a binding constraint instead of assuming that it is one.

\textbf{The policy implications are substantial. If domestic distributive conflict is the main obstacle to climate action, then institutions designed primarily to monitor and punish international defection may target the wrong problem}. International cooperation may be more effective when it empowers pro-climate constituencies, weakens or compensates politically powerful losers, supports clean-energy coalitions, and changes domestic distributions of political influence. The article therefore calls for a richer research program that examines how domestic and international factors interact rather than treating the global climate problem as a self-evident Prisoner's Dilemma.


\clearpage
\section{Bechtel, Genovese, and Scheve --- \textit{Interests, Norms and Support for the Provision of Global Public Goods: The Case of Climate Co-operation}}

\subsection{Research question and theoretical framework}

Bechtel, Genovese, and Scheve investigate the domestic foundations of public support for international climate cooperation. Mitigating climate change requires states to provide a global public good, but international agreements must ultimately be implemented through domestic politics. Public opinion therefore matters because democratic governments may be more willing to negotiate and enforce costly mitigation policies when citizens support them.

The central question is why some individuals favor international climate cooperation while others oppose it. The authors connect this problem to a broader social-science debate between explanations based on \textit{interests} and explanations based on \textit{norms}. The first corresponds to a logic of consequences: individuals support policies that maximize their material welfare. The second corresponds to a logic of appropriateness: individuals support policies that fit internalized ideas about appropriate and socially desirable behavior.

International mitigation resembles a public-goods problem in which norms of reciprocity and altruism can sustain cooperation. The contribution of the article is therefore to show that economic interests and social norms are complementary foundations of support for global climate governance.

\subsection{Sectoral economic interests}

The authors therefore predict that individuals working in more pollution-intensive sectors will be less supportive of international climate agreements. Their empirical strategy improves on subjective measures of economic self-interest by linking each respondent's sector of employment to objective industry-level data on greenhouse-gas emissions and other indicators of pollution intensity.

The argument is explicitly distributive. Even when a society as a whole benefits from avoiding dangerous climate change, the domestic costs of mitigation are concentrated.

\subsection{Social norms: reciprocity and altruism}

The second part of the theory focuses on two other-regarding social norms.:

\textit{Reciprocity} is the willingness to cooperate when others also cooperate and to react negatively to uncooperative behavior. Applied to climate politics, citizens with stronger reciprocal tendencies should be more supportive of international agreements in which other countries also participate.

\textit{Altruism} is an unconditional concern for the welfare of others. Climate mitigation produces benefits for people in other countries and for future generations. Individuals who attach greater weight to others' welfare should therefore be more willing to accept the costs of climate action even when the benefits are diffuse and long term.

\subsection{Data and measurement}

The study uses original large-scale surveys fielded in 2012 in France, Germany, the United Kingdom, and the United States. These countries are wealthy democracies with long histories of greenhouse-gas emissions and are therefore particularly important for international mitigation efforts.

The main dependent variable measures whether respondents support international cooperation to reduce greenhouse-gas emissions even when cooperation involves significant costs.

Authors identify each employed respondent's economic sector and match it with objective pollution data.

Reciprocity is measured through a public-goods game involving potential Amazon gift-card winnings. Respondents specify how much they would give another participant depending on what that person gives them. Individuals whose contribution is more responsive to the partner's contribution are classified as more reciprocal.

Altruism is measured with a separate donation task. Respondents are told that they may win a monetary voucher and can donate part of it to charity. Those willing to make a positive donation are classified as relatively altruistic. Importantly, the reciprocity and altruism measures are only weakly correlated, supporting the claim that they capture distinct normative orientations.

\subsection{Observational results}

Individuals employed in high-emission sectors are significantly less supportive of international climate cooperation.

Social norms: Individuals with above-median reciprocity scores are about 10 percentage points more likely to support global climate policy, and highly altruistic individuals are also roughly 10 percentage points more supportive.

The evidence does not support a simple either/or model in which material self-interest crowds out normative considerations or vice versa. Instead, both help explain the same individual's policy preferences.

\subsection{Conjoint experiment}

Respondents choose between hypothetical international climate agreements whose characteristics vary randomly. The design allows the authors to identify the causal effect of particular treaty features on public support.

The experiment confirms that \textbf{citizens care strongly about material costs}. Increasing the household cost of an agreement from the lowest to higher levels substantially reduces support. Moving from a cost equivalent to 0.5 percent of GDP to one equivalent to 1 percent of GDP lowers support by about 10 percentage points. Respondents also react to sanctions: small sanctions can be acceptable, but medium and high sanctions reduce support.

\textbf{Sectoral interests shape these reactions}. Individuals working in high-emission industries are especially sensitive to sanctions, plausibly because they expect their sectors to face greater difficulty meeting ambitious targets and therefore perceive a higher risk of bearing the associated costs. A high sanction has a substantially larger negative effect among workers in more polluting sectors than among workers in cleaner sectors.

The experiment also shows that \textbf{international participation matters}. Agreements involving more countries or a larger share of global emissions receive more support. This is consistent with \textbf{reciprocity} and conditional cooperation. However, one of the article's important findings is that material interests can moderate this normative response: workers in pollution-intensive sectors are less responsive to broader participation than workers in cleaner sectors. For them, the expected costs of climate regulation remain highly salient even when other countries join.

\subsection{Implications}

The article shows that domestic conflict over international climate cooperation cannot be reduced either to material self-interest or to normative values. Distributional exposure matters because mitigation imposes unequal sectoral costs. Norms matter because some citizens are more willing than others to contribute to collective goods, reciprocate cooperation, and protect the welfare of distant or future beneficiaries.

For policy design, this implies two complementary strategies. Governments can increase support by managing the distributional consequences of climate policy, for example by spreading costs more broadly, compensating vulnerable groups, or facilitating adjustment in pollution-intensive sectors. At the same time, institutional design can make use of other-regarding preferences by increasing visible participation, emphasizing reciprocity, and structuring cooperation in ways that are consistent with citizens' fairness and altruistic concerns.

More broadly, the study contributes to debates in international relations by demonstrating at the individual level that the logic of consequences and the logic of appropriateness can operate together. Interests and norms are not competing master explanations but interacting foundations of political preferences.


\clearpage
\section{Dietz, Ostrom, and Stern --- \textit{The Struggle to Govern the Commons}}

\subsection{Beyond the ``tragedy of the commons''}

Dietz, Ostrom, and Stern revisit the problem of governing common-pool resources and reject the simple conclusion that environmental commons can only be protected through centralized state control or privatization. Hardin's classic ``tragedy of the commons'' correctly drew attention to growing pressure on natural resources, but it underestimated the capacity of communities to create effective self-governing institutions. Empirical research shows successes and failures under community, private, and government ownership alike; no single ownership form guarantees sustainability.

Governance is especially difficult because social and ecological conditions change. Institutions that work under one set of conditions can become ineffective as technologies, markets, populations, or ecosystems evolve. Effective governance must therefore be adaptive rather than fixed.

The authors identify conditions under which commons governance is easier: resource conditions and users' behavior can be monitored at reasonable cost; ecological and social change is moderate; users communicate frequently and possess dense social networks that sustain trust; outsiders can be excluded from resource use; and users regard monitoring and enforcement as legitimate. These conditions are increasingly difficult to obtain as environmental problems become larger, more global, and more tightly connected to external markets.

\subsection{Requirements for adaptive governance}

The central governance challenge is to build institutions capable of operating under uncertainty, conflict, and scale mismatch.

\begin{enumerate}
    \item \textbf{Reliable Information:} Governance systems need knowledge about resource stocks, flows, human behavior, uncertainty, and relevant social values. Information must match the scale of the problem and be trusted by users and decision makers. The collapse of the northern cod fishery illustrates the danger of ignoring local knowledge and relying excessively on aggregated scientific estimates.
    
    \item \textbf{Institutions must manage Conflict: } \textbf{Environmental decisions distribute costs and benefits unevenly, so disagreement is unavoidable. Well-designed participatory and deliberative processes can transform conflict into learning rather than paralysis}.
    
    \item \textbf{Governance must induce Compliance}. Graduated sanctions are often more effective than excessively harsh punishment, while legitimacy and local commitment can make informal enforcement powerful. Command-and-control regulation can work when monitoring is feasible, but market instruments, community rules, voluntary mechanisms, and information disclosure each have different strengths and weaknesses.
    
    \item \textbf{Physical and Institutional Infrastructure} matters because it determines monitoring capacity, access to resources, communication, and the feasibility of particular governance arrangements.
    
    \item \textbf{Institutions must be designed for adaptation and learning} because current knowledge is inevitably incomplete.
\end{enumerate}

\subsection{Adaptive governance strategies}

The authors emphasize three strategies for larger-scale environmental problems.

\begin{enumerate}
    \item \textit{Analytic deliberation} brings scientists, resource users, officials, and affected publics into structured dialogue informed by credible analysis. This can improve information, build trust, and facilitate legitimate decisions.
    
    \item \textit{Nesting} means creating complex, redundant, multilevel institutions rather than relying on a single centralized authority. Environmental processes occur at multiple scales, so governance should connect local, regional, national, and global institutions.
    
    \item Finally, \textit{institutional variety} combines hierarchies, markets, and community self-governance. A diversity of overlapping rules and organizations can improve resilience because each institutional form compensates for weaknesses in others.
\end{enumerate}

\clearpage
\section{Roger, Hale, and Andonova --- \textit{The Comparative Politics of Transnational Climate Governance}}

\subsection{Transnational governance and the domestic politics gap}

Roger, Hale, and Andonova examine the rapid growth of transnational climate governance (TCG): arrangements in which \textbf{substate and nonstate actors} located in multiple countries adhere to common rules intended to steer behavior toward public climate objectives. These arrangements include city networks, private standards, business and NGO initiatives, and hybrid public--private schemes.

Participation in TCG varies dramatically across countries, and these differences cannot be understood without examining domestic institutions, policies, civil society, economic structures, and political incentives.

The article therefore develops a comparative political-economy framework centered on the interaction between transnational and domestic politics. The dependent variable is \textit{participation}: the decision by a substate or nonstate actor to join an initiative or claim adherence to its rules. Domestic contexts shape the incentives, opportunities, capacities, and strategies that make such participation more or less likely.

\subsection{Mapping transnational climate governance}

The descriptive evidence shows extraordinary growth. TCG was limited in the early 1990s, expanded slowly during the first years of the UN climate regime, and accelerated sharply after the Kyoto Protocol.

TCG also became more diverse over time. Early initiatives focused heavily on information sharing and networking. Later, standards and commitment schemes expanded rapidly and became the largest category. Operational and financial initiatives also grew. By the period examined, standards and commitments account for roughly 44 percent of initiatives, information and networking for 31 percent, operational schemes for 14 percent, and finance for 11 percent.

The composition of actors changed as well. Early TCG was dominated by public and hybrid arrangements, whereas private participation expanded dramatically after the late 1990s. Entrepreneurial or private schemes eventually became the largest category, accounting for about 46 percent, with hybrid schemes around 40 percent and purely public schemes about 14 percent. Thus, the overwhelming majority of TCG initiatives involve at least one nonstate actor.

At the same time, participation is geographically very uneven. Some countries host large numbers of participants while others host few, and the distribution changes depending on whether the initiative is public, private, or hybrid. This cross-national variation is precisely what motivates the authors' comparative framework.

\subsection{Domestic sources of preferences and strategies}

Authors argue that domestic incentives must be analyzed simultaneously.

Subnational actors may join TCG because they face domestic political pressure from local constituencies or because transnational cooperation helps them achieve local goals such as cleaner air, energy efficiency, sustainable transport, or urban greening.

Evidence from Eastern European cities shows that strong local civil society is an important driver of participation. Transnational engagement is often rooted in domestic constituencies.

NGOs and firms face limited resources and must decide whether to invest in domestic lobbying or in transnational initiatives. Cao and Ward theorize this allocation problem: when domestic political barriers make national regulation difficult, actors may redirect resources toward TCG; when domestic policy is already strong, they may use transnational governance to extend influence abroad. This illustrates why the relationship between domestic and transnational politics is not one-directional.

\subsection{Domestic institutions as opportunity structures}

The second theme is the mediating role of domestic institutions. Transnational actors do not operate in an institutional vacuum. Political systems that guarantee greater civil and political freedoms provide more opportunities for NGOs, firms, cities, and other actors to organize and participate in cross-border governance.

Decentralization also matters. Subnational governments need sufficient legal authority, fiscal autonomy, and administrative capacity to participate effectively in international networks.

Civil society is a particularly important mechanism. Strong local NGOs can provide information, expertise, legitimacy, political pressure, and organizational capacity. The special issue's findings suggest that domestic civil society may sometimes matter more than international NGOs for explaining why cities participate in TCG. Other contributions show that NGOs also connect domestic and transnational arenas by creating governance initiatives and persuading governments to recognize private standards.

The broader implication is that the effect of globalization is filtered through domestic institutions. Similar international pressures can produce different outcomes because countries differ in openness, regulatory structures, decentralization, and the organizational capacity of societal actors.

\subsection{TCG and public policy: substitutes or complements?}

A central debate concerns whether transnational governance compensates for weak state action. TCG grew rapidly during periods of interstate gridlock, especially after the US withdrawal from Kyoto, which encouraged an image of transnational governance as a substitute for failed public regulation.

Countries with stronger climate policies often exhibit more, not less, participation in TCG. Progressive public regulation can increase actors' capacity and incentives to join transnational networks, create markets for climate-friendly technologies, generate domestic expertise, and provide regulatory models that transnational initiatives can extend internationally.

States and international organizations also actively \textit{orchestrate} TCG. A substantial share of initiatives has been created, supported, financed, or otherwise steered by public institutions. The boundary between public and private governance is therefore porous rather than absolute.

TCG can influence domestic policy by supplying expertise, resources, legitimacy, and standards. Green's work on carbon markets shows that national governments sometimes incorporate technical standards developed through transnational governance into formal law.

Many initiatives are not designed to deliver large emission reductions directly. TCG should not be viewed as a replacement for effective national regulation. Its principal importance may lie in complementing, supporting, and reshaping public governance.

\subsection{General lessons and post-Paris implications}

The article identifies several conclusions with relevance beyond climate change.

First, transnational governance should be studied comparatively rather than only through isolated case studies. Large-scale participation data make it possible to identify systematic cross-national variation and distinguish actor-level from country-level causes.

Second, domestic institutions matter: Liberal political systems, decentralization, civil-society capacity, and regulatory structures affect whether substate and nonstate actors can participate internationally. This produces a deeper interdependence between domestic and global governance than accounts focused only on transnational networks would suggest.

Third, micro-level agency remains crucial: Organizational champions, local expertise, administrative capacity, and civil-society resources determine whether opportunities for transnational action are actually used.

Fourth, public and transnational governance are  complementary: Strong domestic policy can generate more transnational activity, while TCG can reinforce domestic regulation. This finding complicates the common assumption that private or transnational arrangements emerge primarily where states fail.

Capacity building is also essential. Cities and firms require technical expertise, internal policy champions, organizational resources, and connections to networks. Governments, donors, and civil society can strengthen TCG by investing in these capacities rather than only launching new international initiatives.

National governments remain central because they structure the opportunity environment in which transnational action occurs. Even governments unwilling to enact comprehensive climate reforms can enable subnational and nonstate participation by providing resources, convening actors, granting regulatory autonomy, and linking domestic organizations to global networks.

\clearpage
\section{Bueno de Mesquita --- \textit{Global Warming: Designing a Solution}}

\subsection{Central thesis}

Climate change is a global problem, but this does not necessarily imply that it requires a global agreement. The benefits of reducing emissions are diffuse and collective, whereas the economic and political costs are immediate and concentrated. Each government therefore prefers other countries to bear those costs. As a result, global agreements tend either to remain shallow or to fail, while smaller agreements, especially bilateral ones, may balance costs and benefits more effectively.

The chapter's central causal chain can be summarized as follows:

\begin{center}
Emissions $\longrightarrow$ global negative externalities $\longrightarrow$ need for cooperation
$\longrightarrow$ incentives to free-ride $\longrightarrow$ shallow agreements
$\longrightarrow$ high formal compliance but limited actual emission reductions.
\end{center}

\subsection{Climate change as a collective-action problem}

\subsubsection{Nature of the problem}

\begin{itemize}
    \item The author describes global warming primarily as an increase in climatic variance and uncertainty.
    \item Greenhouse gases emitted in one country spread globally.
    \item Emissions are therefore a \textit{negative externality}: emitters do not bear the full costs they impose on others.
    \item Mitigation produces a \textit{public good}: everyone benefits from a more stable climate, and potential beneficiaries cannot easily be excluded.
    \item Every actor consequently has an incentive to let others pay for mitigation while continuing to enjoy its benefits.
\end{itemize}

\subsubsection{Unequal distribution of climate risks}

The main threats identified in the chapter are drought, flooding, storms, coastal inundation caused by rising sea levels, and agricultural losses. The countries most exposed to these risks are often poor or middle-income states and generally do not coincide with the largest greenhouse-gas emitters. Climate politics therefore separates two broad groups:

\begin{itemize}
    \item \textbf{Victims}, which demand protection, assistance, or compensation;
    \item \textbf{Major emitters}, which are expected to bear the costs of reducing emissions.
\end{itemize}

\subsection{Collective benefits and individual costs}

In the chapter's model, $B$ represents the collective benefit produced by climate mitigation, while $c$ represents the cost that each participant must bear to reduce its emissions. Every country would prefer to obtain $B$ without paying $c$, thereby creating a strong incentive to free-ride.

Meaningful mitigation requires interventions in agriculture and fertilizers, transport, electricity generation, heating, industry, and the use of fossil fuels. Such policies may produce higher food, fuel, and energy prices; lower production and employment; higher taxation; or cuts to other public programs. Citizens may support climate protection in the abstract while opposing the concrete costs required to achieve it.

\subsubsection{Economic estimates discussed in the chapter}

\begin{itemize}
    \item A McKinsey study estimates a marginal cost below \$50 per ton for some forms of abatement.
    \item Reducing US emissions by approximately 4.5 gigatons would cost around \$225 billion per year, equivalent to roughly 1.5 percent of US GDP at the time.
    \item The National Academies estimate that reductions exceeding 80 percent would be required to stabilize atmospheric carbon-dioxide concentrations.
\end{itemize}

The fundamental question is therefore not only whether mitigation is economically possible, but whether its costs are politically acceptable. Electoral incentives make leaders reluctant to raise taxes, increase prices, eliminate subsidies, or cut popular programs in exchange for benefits that are global and long term.

\subsection{Conflict over the allocation of abatement costs}

\subsubsection{Position of developing countries}

China, India, and other developing countries argue that the problem was largely created by two centuries of Western industrialization. Developed countries therefore bear greater historical responsibility and should assume most of the costs. This position is expressed through the principle of \textit{common but differentiated responsibilities}.

\subsubsection{Position of developed countries}

The United States and European countries emphasize that most future growth in emissions will come from emerging economies. They argue that the problem cannot be solved without binding commitments from all major emitters and that financial assistance must be accompanied by real mitigation, monitoring, and transparency.

The two sides thus rely on different distributive principles:

\begin{itemize}
    \item developed countries look primarily at future emissions;
    \item emerging countries look primarily at historical responsibility;
    \item developed countries demand commitments from all major emitters;
    \item emerging countries demand that wealthy countries bear most of the costs;
    \item developed countries demand monitoring and transparency;
    \item emerging countries seek to preserve economic growth and political autonomy.
\end{itemize}

Both positions are consistent with the domestic political interests of the governments advancing them.

\subsection{The Kyoto Protocol}

\subsubsection{Institutional structure}

The Kyoto Protocol was adopted in 1997 and generally used 1990 emissions as its benchmark. Binding targets applied primarily to developed countries listed in Annex I, while many developing countries could comply without reducing their emissions. The United States signed the Protocol but did not ratify it.

\subsubsection{Why Kyoto was a shallow agreement}

Kyoto obtained broad participation because it demanded little from most signatories. Monitoring depended heavily on national self-reporting, enforcement mechanisms were weak, and states could purchase unused emission allowances from other countries. Economic contraction also reduced emissions in some countries, allowing them to appear successful without adopting new climate policies.

The flexibility mechanism produced mixed effects. It lowered the cost of nominal compliance and transferred resources to poorer economies, but it also allowed states to satisfy accounting targets without necessarily reducing total global emissions. Kyoto therefore achieved high formal compliance while global emissions continued to rise.

\subsection{The failure of Copenhagen}

The 2009 Copenhagen Conference was intended to overcome Kyoto's limitations. The United States proposed mobilizing \$100 billion annually by 2020, subject to three conditions:

\begin{enumerate}
    \item a strong agreement;
    \item meaningful mitigation measures;
    \item transparency in implementation.
\end{enumerate}

These conditions would have created a \textit{deep agreement}, requiring real and costly changes. Copenhagen failed because the proposed funds would have been insufficient if divided among a very large number of countries, China and India rejected the proposed distribution of responsibilities, and China would not accept meaningful international monitoring. Without transparency, violations could neither be reliably detected nor punished.

The contrast between Kyoto and Copenhagen illustrates a general pattern: shallow global agreements attract broad participation, whereas deep global agreements generate resistance and are likely to fail.

\subsection{An illustrative bilateral solution: the United States and Brazil}

The author proposes replacing the pursuit of an unattainable perfect agreement with an achievable \textit{second-best}. A bilateral agreement between the United States and Brazil serves as the principal illustration.

\subsubsection{Why Brazil?}

Brazil is a major emitter, but its population is much smaller than China's. A concentrated financial transfer would therefore generate a larger per-capita benefit and a stronger domestic incentive to comply. According to the chapter's figures, a \$50 billion transfer would amount to approximately \$250 per Brazilian but only about \$37 per Chinese citizen. Moreover, both Brazil and the United States are democracies with large winning coalitions, so the benefits of cooperation would have to reach a substantial share of their populations.

\subsubsection{Proposed structure of the agreement}

\begin{itemize}
    \item The United States provides \$50 billion per year.
    \item Brazil implements agreed and verifiable emission reductions.
    \item Payment is not made in advance.
    \item Funds are held in escrow by a neutral third party and released only after compliance has been verified.
    \item The interaction is repeated annually.
    \item If either party defects, the other responds by withdrawing cooperation, following a \textit{tit-for-tat} strategy.
\end{itemize}

\subsubsection{The repeated Prisoner's Dilemma}

In a one-shot interaction, both governments have an incentive to defect. Repetition, however, creates a \textit{shadow of the future}. Cooperation generates continuing benefits, whereas cheating produces only an immediate advantage followed by the loss of future cooperation. Monitoring, conditional transfers, and credible retaliation can therefore make a small but deep agreement sustainable.

\subsection{Comparison of cooperation models}

\begin{description}
    \item[Shallow global agreement.] Many participants, limited obligations, and weak enforcement produce high participation but few substantive results.
    \item[Deep global agreement.] Costly reductions, effective monitoring, and real sanctions produce strong political resistance and make agreement unlikely.
    \item[Deep bilateral agreement.] Concentrated costs and benefits, targeted transfers, monitoring, and repeated interaction increase the probability of effective cooperation.
    \item[Unilateral action.] A wealthy country may reduce emissions independently, but doing so remains economically and politically costly.
\end{description}

\subsection{Conclusion}

The proposed US--Brazil agreement is not presented as a complete solution to climate change. It illustrates a more politically realistic strategy: concentrate resources on a limited number of major emitters, compensate actors that bear the costs, require measurable reductions, ensure monitoring and punishment, and build cooperation through repeated interaction.

The chapter ultimately argues that small, deep international agreements or credible unilateral initiatives may produce more meaningful results than the continued pursuit of broad global agreements dominated by rhetoric and shallow commitments.

\medskip
\noindent\textit{Note: The economic figures, emission shares, and political examples reflect the book's publication period (2013).}
